% Part: many-valued-logic
% Chapter: three-valued-logics
% Section: goedel

\documentclass[../../../include/open-logic-section]{subfiles}

\begin{document}

\olfileid[pt]{mvl}{thr}{god}

\olsection{Lógicas de G\"odel}

Kurt G\"odel introduziu uma sequência de lógicas $n$-valoradas, cada uma das quais
contém todas as !!{formula}s válidas na lógica intuicionista e está
contida na lógica clássica. Eis a primeira que apresenta interesse:

\begin{defn}\ollabel{defn:goedel}
A \emph{lógica $3$-valorada de G\"odel}~$\LogGod$ é definida pela matriz:
\begin{enumerate}
  \item A linguagem proposicional padrão $\Lang L_0$, com
  $\lfalse$, $\lnot$, $\land$, $\lor$, $\lif$.
  \item O conjunto de valores de verdade $V = \{\True, \Undef, \False\}$.
  \item $\True$ é o único valor designado, isto é, $V^+ = \{\True\}$.
  \item Para $\lfalse$, temos $\tf{\lfalse} = \False$. As funções de
  verdade dos demais conectivos são dadas pelas seguintes
  tabelas:
  \begin{center}
    \begin{tabular}{c|c} 
      $\tf{\lnot}[\LogGod]$ & \\ 
      \hline  
      $\True$ & $\False$ \\ 
      $\Undef$ & $\False$ \\
      $\False$ & $\True$ 
    \end{tabular}
    \quad
    \begin{tabular}{c|ccc} 
      $\tf{\land}[\LogGod]$ & $\True$ & $\Undef$ & $\False$ \\ 
      \hline 
      $\True$ & $\True$ & $\Undef$ & $\False$ \\ 
      $\Undef$ & $\Undef$ & $\Undef$ & $\False$\\ 
      $\False$ & $\False$ & $\False$ & $\False$ 
    \end{tabular}
    \\[2ex]
    \begin{tabular}{c|ccc} 
      $\tf{\lor}[\LogGod]$ & $\True$ & $\Undef$ & $\False$ \\ 
      \hline 
      $\True$ & $\True$ & $\True$ & $\True$ \\ 
      $\Undef$ & $\True$ & $\Undef$ & $\Undef$ \\
      $\False$ & $\True$ & $\Undef$ & $\False$ 
    \end{tabular}
    \quad
    \begin{tabular}{c|ccc} 
      $\tf{\lif}[\LogGod]$ & $\True$ & $\Undef$ & $\False$ \\ 
      \hline 
      $\True$ & $\True$ & $\Undef$ & $\False$ \\ 
      $\Undef$ & $\True$ & $\True$ & $\False$  \\ 
      $\False$ & $\True$ & $\True$ & $\True$ 
    \end{tabular}
  \end{center} 
\end{enumerate}
\end{defn}

Você observará que as tabelas-verdade de $\land$ e~$\lor$ são as
mesmas das lógicas de \L ukasiewicz e forte de Kleene, mas as tabelas-verdade
de $\lnot$ e~$\lif$ diferem em cada uma. Na lógica de G\"odel,
$\tf{\lnot}(\Undef) = \False$. Em contraste com as lógicas de \L ukasiewicz e
Kleene, $\tf{\lif}(\Undef, \False) = \False$; em contraste com a
lógica de Kleene (mas como na de \L ukasiewicz), $\tf{\lif}(\Undef,
\Undef) = \True$.

Como sugere a relação com a lógica intuicionista mencionada acima,
$\LogGod[3]$ é próxima da lógica intuicionista. Todas as verdades intuicionistas
são tautologias em~$\LogGod[3]$, e muitas tautologias clássicas
que não são válidas intuicionisticamente também deixam de ser tautologias
em~$\LogGod[3]$. Por exemplo, as seguintes não são tautologias:
\begin{align*}
  & p \lor \lnot p && (p \lif q) \lif (\lnot p \lor q) \\
  & \lnot\lnot p \lif p && \lnot(\lnot p \land \lnot q) \lif (p \lor q) \\
  & ((p \lif q) \lif p) \lif p && \lnot(p \lif q) \lif (p \land \lnot q)
\end{align*}
Entretanto, nem toda tautologia de $\LogGod[3]$ também é
intuicionisticamente válida; por exemplo, $\lnot\lnot p \lor \lnot p$ ou $(p
\lif q) \lor (q \lif p)$.

\begin{prob}
  Dê tabelas-verdade que mostrem que as seguintes fórmulas são tautologias
  de~$\LogGod[3]$:
  \begin{align*}
    & \lnot\lnot p \lor \lnot p\\
    & (p \lif q) \lor (q \lif p) \\
    & \lnot(p \land q) \lif (\lnot p \lor \lnot q) \\
    & (p \lif q) \lor (q \lif r) \lor (r \lif s)
  \end{align*}
\end{prob}

\begin{prob}
  Dê tabelas-verdade que mostrem que as seguintes fórmulas não são tautologias
  de~$\LogGod[3]$:
  \begin{align*}
    & (p \lif q) \lif (\lnot p \lor q) \\
    & \lnot(\lnot p \land \lnot q) \lif (p \lor q) \\
    & ((p \lif q) \lif p) \lif p \\
    & \lnot(p \lif q) \lif (p \land \lnot q)
  \end{align*}
\end{prob}

\begin{prob}
  Quais das relações a seguir valem na lógica de G\"odel? Dê uma tabela-verdade para cada uma.
  \begin{enumerate}
    \item $p, p \lif q \Entails q$
    \item $p \lor q, \lnot p \Entails q$
    \item $p \land q \Entails p$
    \item $p \Entails p \land p$
    \item $p \Entails p \lor q$
  \end{enumerate}
\end{prob}

\end{document}
